\documentclass[12pt]{article}
\usepackage[utf8]{inputenc}
\usepackage[T1]{fontenc}
\usepackage{hyperref}
\title{Kangri Language Resources for K-12 Curriculum Development}
\author{Compiled by Claude Sonnet 5.5 (claude-sonnet-5-5) for Andr\'as Kornai}
\date{}
\begin{document}
\maketitle

\section{Introduction}
Kangri (ISO 639-3 xnr; Glottocode kang1280) is a Western Pahari language of Himachal Pradesh, India. It is spoken mainly in Kangra district, Hamirpur district and the hilly parts of Una district. The 2011 census gives about 1.1 million speakers, according to Wikipedia. It is written in Takri and Devanagari. UNESCO classifies it as definitely endangered, and it has no official status. It is not known to be taught in schools. The search found no school materials in Kangri, so everything below is linguistic or reference material.

\section{Primary Educational Resources}
The search found no Kangri primers, textbooks or school programmes. A curriculum would have to be built from the descriptive materials below.

\section{Linguistic and Reference Materials}
\begin{itemize}
\item Kangri language, Wikipedia overview: \url{https://en.wikipedia.org/wiki/Kangri_language}
\item ``Kangri In Context: An Areal Perspective'' (dissertation). It covers phonology, tone, stress, morphology and syntax of the Palampur variety, compared with Hindi, Punjabi and Dogri: \url{https://tdl-ir.tdl.org/handle/10106/945}
\item A review of Anil Kumar Swadeshi's \emph{A Grammar of Kangri} (LINCOM), reported in the search to cover phonetics, phonology, morphology, syntax, verb agreement and ergativity: \url{https://linguistlist.org/issues/36/3246}
\item Language Survey of India: Himachal Pradesh. It includes a Kangri chapter with a grammatical sketch and sociolinguistic profile: \url{https://beta-old.ruralindiaonline.org/v1/uploads/Language_Survey_of_India_Himachal_Pradesh_9a8d9d3690.pdf}
\item Kangri numerals (Eaton, MPI lingweb): \url{https://lingweb.eva.mpg.de/channumerals/Kangri.htm}
\end{itemize}

\section{Religious and Community Resources}
The search found no Kangri religious or community resources. No scripture translations or community organisations were identified.

\section{Digital and Audio Resources}
\begin{itemize}
\item Glosbe community dictionary and word lists for xnr. Quality is unverified and it is user-built: \url{https://glosbe.com/lib/xnr}
\end{itemize}
No audio or video resources were found in these searches.

\section{Curriculum Development Context}
Kangri has many speakers but almost no educational infrastructure. The search found no standalone printed Kangri--English dictionary, and no primers or readers. A curriculum would need new materials in Takri and Devanagari, using the grammars above as a base. Fieldwork with the community would be needed to fill the gaps.

\end{document}